\documentclass[10pt,a4paper]{amsart}
\usepackage[margin=19mm]{geometry}
\usepackage{amsmath,amssymb,mathtools}
\usepackage[scr=rsfs,cal=euler]{mathalfa}
\usepackage{xfrac}
\usepackage[unicode,hidelinks,bookmarks=false]{hyperref}
\newcommand{\ie}{i.e.\ }
\newcommand{\cf}{cf.\ }
\newcommand{\resp}{respectively\ }
\newcommand{\Subsection}{\subsection}
\providecommand{\nobreakdash}{\nobreak\hbox{-}}
\setlength{\emergencystretch}{2em}
\multlinegap0pt
\usepackage{amssymb}
\usepackage[shortlabels]{enumitem}
\newtheorem{definition}{Definition}
\newtheorem{lemma}{Lemma}
\newtheorem{theorem}{Theorem}
\newcommand{\Dom}{\mathcal{D}} 
\newcommand{\R}{\mathbb{R}}
\DeclareMathOperator{\SO}{SO} 
\DeclareMathOperator{\dbsup}{db\text{-}sup}
\DeclareMathOperator{\tr}{tr}
\title{A computer-assisted proof of Kuperberg's six-cylinder conjecture}
\author{Ivan Mati\'c, Rados Radoi\v{c}i\'c}
\date{Source: arXiv:2607.24691v2 (CC BY 4.0)}
\begin{document}
\maketitle

\noindent\emph{Source and licence.} A computer-assisted proof of Kuperberg's six-cylinder conjecture. Version arXiv:2607.24691v2 (CC BY 4.0), distributed by arXiv under the Creative Commons Attribution 4.0 International licence, http://creativecommons.org/licenses/by/4.0/, as recorded in the arXiv licence metadata for this exact version and shown on the abstract page https://arxiv.org/abs/2607.24691v2. Full licence notice: \texttt{licenses/arxiv-2607.24691-CC-BY-4.0.txt}.\par\smallskip

\noindent\textit{Editorial scope.} Counted unit: the article's theorem thm:dbBound (source0.tex lines 240--248). The article's complete original proof of it is delivered with it (source0.tex lines 249--278, one \texttt{proof} environment of 30 lines). No mathematics is written, repaired or re-derived: the statement and the proof are the article's own lines, in the article's own notation, and no retained line is substituted or re-flowed.  Retained uncounted as the source-local context the proof needs: lines 177--180; lines 205--210.  Nothing was omitted from any retained span, so the package carries no omission record.  No retained line contains a citation, so the wrapper carries no bibliography and no bibliography entry.  No retained line contains a figure environment, an image-inclusion command or drawing code, so no image is delivered and the package declares no asset dependency.  Editorial formatting only, in addition: an AMS \texttt{amsart} preamble that restates the article's own declarations and macros used by the retained lines, namely the environments \texttt{definition}, \texttt{lemma}, \texttt{theorem} on separate counters, together with the standard packages those lines need.  The retained environments print in this document as Definition 1, Lemma 1, Theorem 1 in order of appearance, whereas the article prints the same items with the numbering of its own full document.  The complete native source of the article as distributed in the arXiv e-print for this exact version is bundled unchanged as \texttt{sources/206137/source0.tex}.
\begin{definition}\label{def:nonOverlappingX}
A set $X\subseteq \SO(3)$ will be called \emph{non-overlapping} if
$g^{-1}g'\in\Dom$ whenever $g$ and $g'$ are two distinct elements of $X$.
\end{definition}

\begin{lemma}\label{lem:Torth}
Under (i)--(ii), for every $g\in\SO(3)$: $\det T^\ell(g)=\pm1$; $T^\ell(e)=I$; and
\[
T^\ell(g^{-1})\;=\;T^\ell(g)^{-1}\;=\;D_\ell^{-1}\,T^\ell(g)^\top D_\ell .
\]
\end{lemma}

\begin{theorem}\label{thm:dbBound}
Let $\overrightarrow C=(C_\ell)_{\ell=1}^{L}$ be a sequence of symmetric positive-semidefinite
matrices $C_\ell\in\R^{n_\ell\times n_\ell}$. If $\dbsup(\overrightarrow C)<0$ and $M$ is any real
number with $\dbsup(\overrightarrow C)\le M<0$, then
\begin{equation}\label{eq:dbBound}
N_{\max}\;\le\;1+\left\lfloor
\frac{\sum_{\ell=1}^{L}\tr\bigl(D_\ell^{-1}C_\ell\bigr)}{-M}\right\rfloor .
\end{equation}
\end{theorem}

\begin{proof}
Let $g_1,\dots,g_N$ be distinct elements of a non-overlapping set
(Definition~\ref{def:nonOverlappingX}), and put $S_\ell=\sum_{i=1}^{N}T^\ell(g_i)$. By
Lemma~\ref{lem:Torth},
\[
T^\ell\bigl(g_i^{-1}g_j\bigr)=T^\ell(g_i)^{-1}T^\ell(g_j)
=D_\ell^{-1}\,T^\ell(g_i)^\top D_\ell\,T^\ell(g_j),
\]
so that, summing over all ordered pairs and using linearity and cyclicity of the trace,
\begin{align*}
\sum_{i=1}^{N}\sum_{j=1}^{N}\tr\!\Bigl(D_\ell^{-1}C_\ell\,T^\ell\bigl(g_i^{-1}g_j\bigr)\Bigr)
&=\sum_{i,j}\tr\!\Bigl(D_\ell^{-1}C_\ell D_\ell^{-1}\,T^\ell(g_i)^\top D_\ell\,T^\ell(g_j)\Bigr)\\
&=\tr\!\Bigl(\bigl(D_\ell^{-1}C_\ell D_\ell^{-1}\bigr)\bigl(S_\ell^\top D_\ell S_\ell\bigr)\Bigr)
\;\ge\;0,
\end{align*}
since the trace of a product of two symmetric positive-semidefinite matrices is non-negative
(for such $P,Q$: $\tr(PQ)=\tr(\sqrt Q\,P\sqrt Q\,)\ge0$, the trace of a positive-semidefinite
matrix): $D_\ell^{-1}C_\ell D_\ell^{-1}$ is a congruence of $C_\ell$, and
$S_\ell^\top D_\ell S_\ell$ a congruence of $D_\ell$, so both are symmetric positive semidefinite.

Summing over $\ell$ and separating the $N$ diagonal pairs $i=j$ (where
$T^\ell(e)=I$ gives the value $\sum_\ell\tr(D_\ell^{-1}C_\ell)$) from the $N(N-1)$ off-diagonal
pairs (where $g_i^{-1}g_j\in\Dom$ by Definition~\ref{def:nonOverlappingX}, so each term is at most
$\dbsup(\overrightarrow C)\le M$),
\[
0\;\le\;N\sum_{\ell=1}^{L}\tr\bigl(D_\ell^{-1}C_\ell\bigr)\;+\;N(N-1)\,M .
\]
Since $M<0$, dividing by $N$ and rearranging gives
$(N-1)(-M)\le\sum_\ell\tr(D_\ell^{-1}C_\ell)$, which implies \eqref{eq:dbBound}.
\end{proof}

\end{document}
